\subsection{Counting with Eisenstein norms}\label{sec:norm}
We begin by finding plenty of candidate pairs. A quaternion packages two
Eisenstein norms into one multiplicative norm, which lets us count the
representations exactly. This count is the supply of pairs from which we
will remove the bad ones.

For positive integers, write $\sigma(n)=\sum_{d\mid n}d$ and
$\tau(n)=\sum_{d\mid n}1$. A term $\sigma(n/4)$ is interpreted as zero
when $4\nmid n$. Let $\omega^2+\omega+1=0$, $R=\Z[\omega]$, and
\[
Q(x,y)=x^2-xy+y^2=N_{\Q(\omega)/\Q}(x+y\omega).
\]
For instance, $Q(3,1)=7$ and $Q(4,1)=13$, so $20$ is represented by
$Q(x,y)+Q(u,v)$ with both norms already in $\AA$.

Consider the definite quaternion algebra and order
\[
B=\Q(\omega)\oplus\Q(\omega)j,\qquad
j^2=-1,\quad jz=\bar z j,\qquad \OO=R\oplus Rj.
\]
Its reduced norm is
\[
\nrd(z+w j)=N(z)+N(w).
\]
In the basis $1,\omega,j,\omega j$, the trace pairing has determinant $-9$.
Consequently $\OO$ has reduced discriminant $3$. Definiteness and the parity
of the ramification set imply that the algebra discriminant is $3$, so
$\OO$ is maximal.

Every complex number is within squared distance $1/3$ of an Eisenstein
integer: the equilateral triangles in the lattice have circumradius
$1/\sqrt3$. Rounding the two complex coordinates independently gives
\[
\forall x\in B\otimes\mathbb R\quad\exists q\in\OO:
\quad\nrd(x-q)\le2/3<1.
\]
The one-coordinate rounding bound also makes $R$ norm Euclidean, hence
a unique factorization domain. The two-coordinate bound makes $\OO$
norm Euclidean on either side. In a nonzero integral right
ideal, choose $\alpha$ of least positive norm and divide each $\beta$ by
$\alpha$ on the left. A nonzero remainder would have smaller norm, so the
ideal is $\alpha\OO$. Fractional ideals reduce to this case by scaling;
the right ideal class number is one.

The units are precisely $R^\times\cup R^\times j$, and there are twelve.
For $3\nmid n$ let
\[
\mathcal Q_n=\{\alpha\in\OO:\nrd(\alpha)=n\},\qquad r(n)=|\mathcal Q_n|.
\]
At a prime $\ell\mid n$ the order is locally $M_2(\Z_\ell)$. There are
$1+\ell+\cdots+\ell^e$ integral right ideals of reduced norm $\ell^e$,
by the index-$\ell^e$ sublattice count in $\Z_\ell^2$. Multiplication of
the local counts and principality give $\sigma(n)$ right ideals of norm
$n$, each with twelve generators of that norm. Hence
\begin{equation}\label{eq:count}
r(n)=12\sigma(n)\qquad(3\nmid n).
\end{equation}
This is the usual all-integral-ideal normalization of a Brandt operator;
see \cite[Proposition 41.3.1]{Voight}.

For $n\equiv2\pmod3$, both component norms in every representation
$n=Q(x,y)+Q(u,v)$ are $1\pmod3$, because
$Q(x,y)\equiv(x+y)^2\pmod3$ is either $0$ or $1$. In particular both
summands are positive and neither is divisible by $3$.
